% TODO proof-reading
\subsubsection{\main関数の唯一の引数も配列です}

\main関数の唯一の引数である文字列の配列を使用します：


\begin{lstlisting}[style=customjava]
public class UseArgument
{
	public static void main(String[] args)
	{
		System.out.print("Hi, ");
		System.out.print(args[1]);
		System.out.println(". How are you?");
	}
}
\end{lstlisting}

0番目の引数はプログラム名です（\CCppなどと同様）、
従ってユーザーが提供する1番目の引数は1番です。


\begin{lstlisting}
  public static void main(java.lang.String[]);
    flags: ACC_PUBLIC, ACC_STATIC
    Code:
      stack=3, locals=1, args_size=1
         0: getstatic     #2      // Field java/lang/System.out:Ljava/io/PrintStream;
         3: ldc           #3      // String Hi, 
         5: invokevirtual #4      // Method java/io/PrintStream.print:(Ljava/lang/String;)V
         8: getstatic     #2      // Field java/lang/System.out:Ljava/io/PrintStream;
        11: aload_0       
        12: iconst_1      
        13: aaload        
        14: invokevirtual #4      // Method java/io/PrintStream.print:(Ljava/lang/String;)V
        17: getstatic     #2      // Field java/lang/System.out:Ljava/io/PrintStream;
        20: ldc           #5      // String . How are you?
        22: invokevirtual #6      // Method java/io/PrintStream.println:(Ljava/lang/String;)V
        25: return        
\end{lstlisting}

11の\TT{aload\_0}は0番目の\ac{LVA}スロット（1番目で唯一の\main引数）の\IT{参照}をロードします。

12と13の\TT{iconst\_1}と\TT{aaload}は、配列の最初（0から数える）の要素への\IT{参照}を取得します。

文字列オブジェクトへの\IT{参照}は、オフセット14で\ac{TOS}にあり、そこから\TT{println}メソッドによって取得されます。